\documentclass[11pt,a4paper]{article}
\usepackage[utf8]{inputenc}
\usepackage[T1]{fontenc}
\usepackage{amsmath,amssymb,amsthm,mathtools}
\usepackage[margin=2.5cm]{geometry}
\usepackage{booktabs}
\usepackage{array}
\usepackage{hyperref}
\usepackage{listings}
\lstset{basicstyle=\small\ttfamily,breaklines=true}
\setlength{\emergencystretch}{2em}

\theoremstyle{plain}
\newtheorem{theorem}{Theorem}[section]
\newtheorem{lemma}[theorem]{Lemma}
\newtheorem{corollary}[theorem]{Corollary}
\newtheorem{proposition}[theorem]{Proposition}
\newtheorem{conjecture}[theorem]{Conjecture}
\theoremstyle{definition}
\newtheorem{definition}[theorem]{Definition}
\theoremstyle{remark}
\newtheorem{remark}[theorem]{Remark}

\newcommand{\PCF}{\mathrm{PCF}}
\newcommand{\half}{\tfrac{1}{2}}
\newcommand{\dff}[1]{(#1)!!}
\newcommand{\hyp}[5]{{{}_{#1}F_{#2}}\!\left(\genfrac{}{}{0pt}{}{#3}{#4}\bigg|#5\right)}

\title{Polynomial Continued Fractions: a Proved Logarithmic Ladder,\\
a \texorpdfstring{$4/\pi$}{4/pi} Casoratian Identity, and 482 Irrational Constants}
\author{Papanokechi}
\date{April 2026}

\begin{document}
\maketitle

\begin{abstract}
We prove that the polynomial continued fractions
$\PCF(-kn^2,\,(k{+}1)n{+}k)$ and $\PCF(-n(2n{-}3),\,3n{+}1)$
converge respectively to $1/\ln(k/(k{-}1))$ for every real $k>1$
and to $4/\pi$, and from the second family we derive a rapidly
convergent identity for $\pi/4$.
The key tool is a combination of closed-form convergents and a
discrete Casoratian telescope for the underlying three-term recurrence.
This differs from the classical Gauss and Brouncker formulas in that
the evaluation is obtained directly inside the recurrence, with error
$O(2^{-N}/N^{7/2})$ for the resulting series.
As a computational application, we certify 482 quadratic generalized
continued fractions against a basis of 15 standard constants using Arb
ball arithmetic at 10,000-bit precision and depth $N=5000$.
These computations yield $1500{+}$ correct digits and rigorous
irrationality statements for 482 previously unclassified constants.
We also prove that the $4/\pi$ family cannot be expressed as a Gauss
${}_2F_1$ continued fraction, because its term ratio is forced into a
${}_2F_3$ form over $\mathbb{Q}(\sqrt{5})$.
\end{abstract}

\noindent
\textbf{2020 Mathematics Subject Classification:}
11J70 (primary), 11Y60, 33C20.

\noindent
\textbf{Keywords:}
continued fractions, irrationality, Casoratian,
Gauss hypergeometric, Pincherle's theorem, PSLQ.

\tableofcontents

%======================================================================
\section{Introduction}\label{sec:intro}
%======================================================================

Two explicit one-parameter families of polynomial continued fractions admit
closed evaluations here. One family produces a full logarithmic ladder
$1/\ln(k/(k{-}1))$ for every real $k>1$; another yields $4/\pi$, a rapidly
convergent series for $\pi/4$, and an explicit obstruction showing that the
same continued fraction is not of Gauss ${}_2F_1$ type.

Classical work of Wall~\cite{wall1948} and Lorentzen--Waadeland~\cite{lorentzen2008}
provides the analytic background on continued fractions, convergence criteria,
and the Gauss hypergeometric model relevant here. Ap\'ery's use of a polynomial
continued fraction in the irrationality proof for $\zeta(3)$~\cite{apery1979}
shows that such recurrences can carry arithmetic content, while the Ramanujan
Machine program~\cite{raayoni2021,elimelech2024} has renewed the search for
low-degree examples with closed values.

A \emph{polynomial continued fraction} (PCF) with numerator polynomial
$\alpha(n)$ and denominator polynomial~$\beta(n)$ is
\begin{equation}\label{eq:pcf}
  v = \beta(0) + \cfrac{\alpha(1)}{\beta(1) +
  \cfrac{\alpha(2)}{\beta(2) + \cfrac{\alpha(3)}{\beta(3) + \cdots}}}.
\end{equation}
The convergents $C_n = p_n/q_n$ are computed via the forward recurrence
\begin{equation}\label{eq:rec}
  p_n = \beta(n)\,p_{n-1} + \alpha(n)\,p_{n-2}, \qquad
  q_n = \beta(n)\,q_{n-1} + \alpha(n)\,q_{n-2},
\end{equation}
with $p_{-1} = 1$, $p_0 = \beta(0)$, $q_{-1} = 0$, $q_0 = 1$.

The proved statements are separated carefully from the conjectural ones.
Section~\ref{sec:log} establishes the logarithmic ladder for all real $k>1$.
Sections~\ref{sec:pi} and~\ref{sec:closed} prove the base cases $m=0$ and
$m=1$ of the $\pi$-family, and Section~\ref{sec:casoratian} derives the
Casoratian identity for $\pi/4$. Section~\ref{sec:gauss} proves non-membership
in the Gauss continued-fraction class, and Section~\ref{sec:vquad} proves
irrationality for the quadratic catalogue. What remains conjectural is explicit:
Conjecture~\ref{conj:ratio} proposes the general ratio law in the parameter
$m$, and Section~\ref{sec:gamma} derives the continuous Gamma quotient from
that conjecture together with the proved base cases. Sections~\ref{sec:euler}--\ref{sec:meta}
place the proved formulas in context, while the appendices record numerical
certification, the search procedure, exact convergents, and reproducibility data.

%======================================================================
\section{The Logarithmic Ladder}\label{sec:log}
%======================================================================

\begin{theorem}[Logarithmic Ladder]\label{thm:log}
For all real $k > 1$, the PCF with
$\alpha(n) = -kn^2$ and $\beta(n) = (k{+}1)n + k$
converges to $1/\ln(k/(k{-}1))$.
\end{theorem}

\begin{proof}
\textbf{Claim.}
\begin{equation}\label{eq:log-pq}
  p_n = (n{+}1)!\, k^{n+1}, \qquad
  q_n = (n{+}1)!\, \sum_{j=0}^{n} \frac{k^{n-j}}{j+1}.
\end{equation}

\textbf{Base case.}
$p_0 = \beta(0) = k = 1!\cdot k^1$ and $q_0 = 1 = 1!\cdot k^0/1$.

\medskip\noindent\textbf{Induction for $p_n$.}
Assume $p_{n-1} = n!\,k^n$ and $p_{n-2} = (n{-}1)!\,k^{n-1}$.
\begin{align*}
  \beta(n)\,p_{n-1} + \alpha(n)\,p_{n-2}
  &= \bigl((k{+}1)n + k\bigr) n!\,k^n - kn^2 (n{-}1)!\,k^{n-1} \\
  &= n!\,k^n\bigl((k{+}1)n + k - n\bigr)
   = n!\,k^n \cdot k(n{+}1) = (n{+}1)!\,k^{n+1}.
\end{align*}

\medskip\noindent\textbf{Induction for $q_n$.}
Write $S_n = \sum_{j=0}^{n} k^{n-j}/(j{+}1)$
and note $S_n = k\,S_{n-1} + 1/(n{+}1)$.
Assume $q_{n-1} = n!\,S_{n-1}$ and $q_{n-2} = (n{-}1)!\,S_{n-2}$.
Using $S_{n-1} = k\,S_{n-2} + 1/n$:
\begin{align*}
\beta(n)\,q_{n-1} + \alpha(n)\,q_{n-2}
&= \bigl((k{+}1)n{+}k\bigr)\,n!\bigl(k\,S_{n-2}+\tfrac{1}{n}\bigr)
   - kn^2\,(n{-}1)!\,S_{n-2} \\
&= S_{n-2}\cdot k^2(n{+}1)\,n!
   + (n{-}1)!\bigl(kn+k+n\bigr).
\end{align*}
Meanwhile,
$(n{+}1)!\,S_n
= k^2(n{+}1)\,n!\,S_{n-2} + k(n{+}1)(n{-}1)! + n!$.
Since $(n{-}1)!(kn{+}k{+}n) = k(n{+}1)(n{-}1)! + n!$,
the two expressions are equal.

\medskip\noindent\textbf{Limit.}
The closed forms~\eqref{eq:log-pq} are polynomial identities in~$k$,
valid for all real~$k$.
The convergent ratio is
$C_n = k/\sum_{j=0}^{n}(1/k)^j/(j{+}1)$.
For $|1/k| < 1$, the series
$\sum_{j=0}^{\infty} x^j/(j{+}1) = -\ln(1{-}x)/x$
converges at $x = 1/k$, giving
$\sum (1/k)^j/(j{+}1) = k\ln(k/(k{-}1))$.
Therefore
$\lim C_n = k/(k\ln(k/(k{-}1))) = 1/\ln(k/(k{-}1))$.
\end{proof}

\begin{remark}
The theorem holds for all real $k > 1$, not just integers.
Using Arb ball arithmetic at $120$ decimal digits and forward depth
$N=2000$, we checked $15$ non-integer values in $[1.5,100]$ and found
$90$--$120$ matching digits with $1/\ln(k/(k{-}1))$.
Table~\ref{tab:log} lists representative instances.
\end{remark}

\begin{table}[ht]
\centering
\caption{Logarithmic Ladder: first instances}\label{tab:log}
\begin{tabular}{@{}cccc@{}}
\toprule
$k$ & $\alpha(n)$ & $\beta(n)$ & Value \\
\midrule
$2$ & $-2n^2$ & $3n+2$ & $1/\ln 2$ \\
$3$ & $-3n^2$ & $4n+3$ & $1/\ln(3/2)$ \\
$4$ & $-4n^2$ & $5n+4$ & $1/\ln(4/3)$ \\
$5$ & $-5n^2$ & $6n+5$ & $1/\ln(5/4)$ \\
$\varphi$ & $-\varphi\, n^2$ & $(\varphi{+}1)n{+}\varphi$ &
  $1/\ln(\varphi/(\varphi{-}1))$ \\
$k$ & $-kn^2$ & $(k{+}1)n{+}k$ & $1/\ln(k/(k{-}1))$ \\
\bottomrule
\end{tabular}
\end{table}

%======================================================================
\section{The Pi Family}\label{sec:pi}
%======================================================================

\subsection{Base case: \texorpdfstring{$2/\pi$}{2/pi}}\label{sec:pi_base}

\begin{theorem}[Pi Family, $m=0$]\label{thm:pi0}
The PCF with $\alpha(n) = -n(2n{-}1)$ and $\beta(n) = 3n+1$
converges to $2/\pi$.
\end{theorem}

\begin{proof}
\textbf{Claim.}
\begin{equation}\label{eq:pi-pq}
  p_n = (2n{+}1)!!, \qquad
  q_n = (2n{+}1)!! \sum_{j=0}^{n} \frac{j!}{(2j{+}1)!!}.
\end{equation}

\textbf{Base case.}
$p_0 = 1 = (1)!!$ and $q_0 = 1 = 1\cdot(0!/1!!)$.

\medskip\noindent\textbf{Induction for $p_n$.}
Assume $p_{n-1} = (2n{-}1)!!$ and $p_{n-2} = (2n{-}3)!!$.
Using $(2n{-}1)!! = (2n{-}1)(2n{-}3)!!$:
\begin{align*}
  \beta(n)\,p_{n-1} + \alpha(n)\,p_{n-2}
  &= (3n{+}1)(2n{-}1)!! - n(2n{-}1)(2n{-}3)!! \\
  &= (2n{-}3)!!\bigl[(3n{+}1)(2n{-}1) - n(2n{-}1)\bigr] \\
  &= (2n{-}3)!!(2n{-}1)(2n{+}1) = (2n{+}1)!!.
\end{align*}

\medskip\noindent\textbf{Induction for $q_n$.}
Write $T_n = \sum_{j=0}^{n} j!/(2j{+}1)!!$.
Note $T_{n-1} = T_{n-2} + (n{-}1)!/(2n{-}1)!!$.
With inductive hypotheses
$q_{n-1} = (2n{-}1)!!\,T_{n-1}$ and $q_{n-2} = (2n{-}3)!!\,T_{n-2}$:
\begin{align*}
\beta(n)\,q_{n-1} + \alpha(n)\,q_{n-2}
&= (2n{-}1)!!\bigl[(3n{+}1)\,T_{n-1} - n\,T_{n-2}\bigr].
\end{align*}
Substituting $T_{n-1} = T_{n-2} + (n{-}1)!/(2n{-}1)!!$:
\[
  (3n{+}1)\,T_{n-1} - n\,T_{n-2}
  = (2n{+}1)\,T_{n-2} + \frac{(3n{+}1)(n{-}1)!}{(2n{-}1)!!}.
\]
The identity reduces to
$(3n{+}1)(n{-}1)! = (2n{+}1)(n{-}1)! + n!$,
which holds since $n! = n(n{-}1)!$.

\medskip\noindent\textbf{Limit.}
From the closed forms~\eqref{eq:pi-pq}, we have $C_n = p_n/q_n = 1/T_n$.
The identity~\cite[\S5.2]{cuyt2008}
\begin{equation}\label{eq:pi-id}
  \frac{\pi}{2}
  = \sum_{j=0}^{\infty} \frac{j!}{(2j{+}1)!!}
  = \sum_{j=0}^{\infty} \frac{2^j(j!)^2}{(2j{+}1)!},
\end{equation}
derivable from $\arcsin(1) = \pi/2$,
gives $\lim T_n = \pi/2$ and $\lim C_n = 2/\pi$.
\end{proof}

\subsection{Parity phenomenon}\label{sec:parity}

\begin{theorem}[Parity]\label{thm:parity}
For even $c = 2\ell$, the PCF with $\alpha(n) = -n(2n{-}c)$ and
$\beta(n) = 3n+1$ evaluates to an exact rational.
\end{theorem}

\begin{proof}
$\alpha(\ell) = -\ell(2\ell - 2\ell) = 0$, so the CF terminates
at depth~$\ell{-}1$.
\end{proof}

The first even specialisations are listed in Table~\ref{tab:parity}.

\begin{table}[ht]
\centering
\caption{Pi Family: even parameters (exact rationals)}\label{tab:parity}
\begin{tabular}{@{}rcc@{}}
\toprule
$c$ & Value & Formula \\
\midrule
$2$ & $1$ & $1$ \\
$4$ & $3/2$ & $\binom{2}{1}/2$ \\
$6$ & $15/8$ & $\binom{4}{2}/2^3$ \\
$8$ & $35/16$ & $\binom{6}{3}/2^4$ \\
$10$ & $315/128$ & $\binom{8}{4}/2^7$ \\
$12$ & $693/256$ & $\binom{10}{5}/2^8$ \\
\bottomrule
\end{tabular}
\end{table}

\subsection{Odd parameters and the binomial recurrence}

\begin{conjecture}[Ratio Formula]\label{conj:ratio}
Define $\operatorname{val}(m)$ as the limit of the PCF with
$\alpha(n) = -n(2n{-}(2m{+}1))$ and $\beta(n) = 3n+1$.
Then $\operatorname{val}(m{+}1) = \operatorname{val}(m)\cdot 2(m{+}1)/(2m{+}1)$.
\end{conjecture}

\begin{remark}
Using Arb ball arithmetic at $600$ decimal digits and depth $N=5000$,
the recurrence matches to at least $500$ digits for $m = 0,1,\ldots,10$
and for the half-integer values $m = 0.5,1.5,\ldots,10.5$.
It is consistent with the closed form
$\operatorname{val}(m) = 2^{2m+1}/(\pi\binom{2m}{m})$, since
$\binom{2m+2}{m+1} = \binom{2m}{m}\cdot 2(2m{+}1)/(m{+}1)$ gives
ratio $4\binom{2m}{m}/\binom{2m+2}{m+1} = 2(m{+}1)/(2m{+}1)$.
\end{remark}

\begin{corollary}[Full Pi Family]\label{cor:pi}
Assuming Conjecture~\ref{conj:ratio}, for all $m \ge 0$ the PCF
converges to $2^{2m+1}/(\pi\binom{2m}{m})$.
\end{corollary}

\begin{proof}
This is immediate from Theorem~\ref{thm:pi0} and repeated application of
Conjecture~\ref{conj:ratio}. Starting from $\operatorname{val}(0)=2/\pi$,
one obtains
\[
  \operatorname{val}(m)=\frac{2}{\pi}\prod_{j=0}^{m-1}\frac{2(j+1)}{2j+1}
  = \frac{2^{2m+1}}{\pi\binom{2m}{m}},
\]
which is the standard central-binomial product identity.
\end{proof}

Table~\ref{tab:pi} records the first odd specialisations.

\begin{table}[ht]
\centering
\caption{Pi Family: odd parameter $c = 2m{+}1$}\label{tab:pi}
\begin{tabular}{@{}ccccc@{}}
\toprule
$m$ & $c$ & $\alpha(n)$ & $\operatorname{val}(m)\cdot\pi$ &
$\operatorname{val}(m)$ \\
\midrule
$0$ & $1$ & $-n(2n{-}1)$ & $2$ & $2/\pi$ \\
$1$ & $3$ & $-n(2n{-}3)$ & $4$ & $4/\pi$ \\
$2$ & $5$ & $-n(2n{-}5)$ & $16/3$ & $16/(3\pi)$ \\
$3$ & $7$ & $-n(2n{-}7)$ & $32/5$ & $32/(5\pi)$ \\
$4$ & $9$ & $-n(2n{-}9)$ & $256/35$ & $256/(35\pi)$ \\
\bottomrule
\end{tabular}
\end{table}

%======================================================================
\section{Convergent Closed Forms}\label{sec:closed}
%======================================================================

\begin{theorem}[$m=0$ convergent numerators]\label{thm:conv0}
For the Pi Family with $m=0$,
$p_n = (2n{+}1)!! = \dff{2n{-}1}\cdot(2n{+}1)$.
\end{theorem}

\begin{proof}
Base cases: $p_0{=}1$, $p_1{=}3$, $p_2{=}15$.
Induction: assume $p_k = (2k{+}1)!!$ for $k < n$.
The recurrence $p_n = (3n{+}1)p_{n-1} - n(2n{-}1)p_{n-2}$
with $P(n) = 2n{+}1$ reduces to the identity
$(2n{-}1)(2n{-}3)P(n) = (3n{+}1)(2n{-}3)P(n{-}1) - n(2n{-}1)P(n{-}2)$
after dividing by~$\dff{2n{-}3}$.
Expanding both sides yields $0 = 0$.
\end{proof}

\begin{theorem}[$m=1$ convergent numerators]\label{thm:conv1}
For the Pi Family with $m=1$ ($\alpha(n) = -n(2n{-}3)$,
$\beta(n) = 3n{+}1$),
\begin{equation}\label{eq:pn_m1}
  p_n = \dff{2n{-}1}\,(n^2 + 3n + 1).
\end{equation}
The polynomial $h_n = n^2{+}3n{+}1$ (OEIS~\cite{oeis_A028387})
has discriminant~$5$ and
roots $(-3 \pm \sqrt{5})/2$ related to the golden ratio.
\end{theorem}

\begin{proof}
After the substitution $p_n = \dff{2n{-}1}\,u_n$
(using $\dff{2n{-}1} = (2n{-}1)\dff{2n{-}3}$),
the recurrence reduces to
\begin{equation}\label{eq:ured}
  (2n{-}1)\,u_n = (3n{+}1)\,u_{n-1} - n\,u_{n-2}.
\end{equation}
Substituting $u_n = n^2{+}3n{+}1$, $u_{n-1} = n^2{+}n{-}1$,
$u_{n-2} = n^2{-}n{-}1$:
\begin{align*}
  \text{RHS} &= (3n{+}1)(n^2{+}n{-}1) - n(n^2{-}n{-}1) \\
  &= 2n^3 + 5n^2 - n - 1
   = (2n{-}1)(n^2{+}3n{+}1) = \text{LHS.}\qedhere
\end{align*}
\end{proof}

\begin{conjecture}[General $m$]\label{conj:Rnm}
For all $m \ge 0$ and fixed~$n$, the quantity
$R(n,m) = p_n(m)/\dff{2n{-}1}$
is a polynomial in~$m$ of degree $\lfloor n/2 \rfloor$.
\end{conjecture}

A symbolic computation confirms this for $m = 0, \ldots, 12$ and $n \le 20$.
For $m \ge 2$, $R(n,m)$ is rational with denominators dividing the
product of odd primes~$\le 2n{-}1$.

%======================================================================
\section{The Continuous Gamma Identity}\label{sec:gamma}
%======================================================================

The results of \S\S\ref{sec:log}--\ref{sec:closed} are special
cases of a single continuous identity. Its extension beyond the
base cases $m=0,1$ is conditional on Conjecture~\ref{conj:ratio}.

\begin{theorem}[Gamma Function Representation]\label{thm:gamma}
The formula below is proved unconditionally for $m=0$ and $m=1$.
Assuming Conjecture~\ref{conj:ratio}, for all real $m > -\frac{1}{2}$,
the PCF with
$\alpha(n) = -n(2n{-}2m{-}1)$ and $\beta(n) = 3n+1$
converges to
\[
  \operatorname{val}(m)
  = \frac{2\,\Gamma(m{+}1)}{\sqrt{\pi}\,\Gamma(m{+}\half)}.
\]
\end{theorem}

\begin{proof}
We prove the identity for $m = 0$ and $m = 1$ by a purely algebraic
argument, then extend to all~$m$ via the Wallis reduction formula and
Conjecture~\ref{conj:ratio} (supported numerically at $600$ decimal digits and depth $N=5000$).

\medskip
\emph{Step~1: The raising operator.}
Define $L(f)_n = (n{+}2)\,f_n - (n{+}1)^2\,f_{n-1}$.

\begin{lemma}\label{lem:shift}
If\/ $f$ satisfies the $m{=}0$ recurrence
$f_n = (3n{+}1)f_{n-1} - n(2n{-}1)f_{n-2}$,
then $g = L(f)$ satisfies the $m{=}1$ recurrence
$g_n = (3n{+}1)g_{n-1} - n(2n{-}3)g_{n-2}$.
\end{lemma}

\begin{proof}[Proof of Lemma]
Write $g_{n-1} = (n{+}1)f_{n-1} - n^2 f_{n-2}$ and
$g_{n-2} = n\,f_{n-2} - (n{-}1)^2 f_{n-3}$.
Expanding $(3n{+}1)g_{n-1} - n(2n{-}3)g_{n-2}$ and
eliminating $f_{n-3}$ via the $m{=}0$ recurrence at index~$n{-}1$
yields coefficients
\begin{align*}
  \text{coeff.\ of } f_{n-1} &:
    (3n^2{+}4n{+}1) - n(n{-}1) = 2n^2{+}5n{+}1, \\
  \text{coeff.\ of } f_{n-2} &:
    {-}n^2(5n{-}2) + n(n{-}1)(3n{-}2) = {-}n(2n^2{+}3n{-}2),
\end{align*}
matching
$g_n = (n{+}2)[(3n{+}1)f_{n-1} - n(2n{-}1)f_{n-2}]
  - (n{+}1)^2 f_{n-1}$
identically.
\end{proof}

Since $L$ is linear, it maps the solution space~$\mathcal{S}_0$
of the $m{=}0$ recurrence into~$\mathcal{S}_1$.
Checking initial values:
\begin{alignat*}{2}
  L(P^{(0)})_0 &= 2\cdot 1 - 1\cdot 1 = 1 = P_0^{(1)},\;
  &L(P^{(0)})_1 &= 3\cdot 3 - 4\cdot 1 = 5 = P_1^{(1)},\\
  L(Q^{(0)})_0 &= 2\cdot 1 - 1\cdot 0 = 2 = 2\,Q_0^{(1)},\;
  &L(Q^{(0)})_1 &= 3\cdot 4 - 4\cdot 1 = 8 = 2\,Q_1^{(1)}.
\end{alignat*}
By uniqueness:
$P_n^{(1)} = L(P^{(0)})_n$ and $2\,Q_n^{(1)} = L(Q^{(0)})_n$.

\medskip
\emph{Step~2: Poincar\'e growth analysis.}
The recurrence has Poincar\'e characteristic roots $\lambda = 2$
and $\lambda = 1$
(from $\lambda^2 - 3\lambda + 2 = 0$).
Both $P_n^{(m)}$ and~$Q_n^{(m)}$ are dominant, growing
like~$(2n)!!$, so $P_n^{(m)}/Q_n^{(m)} \to S^{(m)}$ at
rate~$O(n!/(2n)!!) = O(2^{-n}/\sqrt{n})$.

\medskip
\emph{Step~3: The limit.}
From Step~1:
\[
  S_n^{(1)}
  = 2\cdot\frac{(n{+}2)P_n^{(0)} - (n{+}1)^2 P_{n-1}^{(0)}}
              {(n{+}2)Q_n^{(0)} - (n{+}1)^2 Q_{n-1}^{(0)}}.
\]
Dividing by $(n{+}2)Q_n^{(0)}$ and using
$(n{+}1)^2/[(n{+}2)(2n{+}1)] \to \frac{1}{2}$,
both correction terms converge to~$\frac{1}{2}$, giving
$S^{(1)} = 2\,S^{(0)}$.

\medskip
\emph{Step~4: General~$m$ and the Gamma identity.}
The base case $S^{(0)} = 2/\pi$ is established by
Theorem~\ref{thm:pi0}.
The Wallis reduction formula gives
$I_{m+1} = \frac{2m{+}1}{2m{+}2}\,I_m$
where $I_m = \int_0^{\pi/2}\sin^{2m}(x)\,dx
= \frac{\sqrt{\pi}\,\Gamma(m{+}\frac{1}{2})}{2\,\Gamma(m{+}1)}$.
Assuming Conjecture~\ref{conj:ratio},
the ratio formula
$\operatorname{val}(m{+}1)/\operatorname{val}(m)
= 2(m{+}1)/(2m{+}1) = I_m/I_{m+1}$
propagates inductively from $S^{(0)} = 1/I_0$,
yielding $\operatorname{val}(m) = 1/I_m =
2\,\Gamma(m{+}1)/(\sqrt{\pi}\,\Gamma(m{+}\half))$
for all non-negative integers~$m$.
Under the same assumption, the identity extends to all real
$m > -\frac{1}{2}$ by analytic continuation.
\end{proof}

\begin{remark}
For $m = 0$ and $m = 1$, Steps~1--3 give a complete algebraic proof
independent of Conjecture~\ref{conj:ratio}:
$S^{(0)} = 2/\pi$ and $S^{(1)} = 4/\pi$.
For $m \ge 2$, the proof is conditional on the ratio conjecture.
No polynomial-coefficient shift operator $L_m$ mapping $\mathcal{S}_m$
to $\mathcal{S}_{m+1}$ has been found for general~$m$~\cite{raayoni2021}.
\end{remark}

\begin{corollary}\label{cor:special}
Assuming Conjecture~\ref{conj:ratio}, the following specializations hold.
\begin{enumerate}
\item[\textup{(i)}]
  Integer $m \ge 0$:
  $\operatorname{val}(m) = 2^{2m+1}/\bigl(\pi\binom{2m}{m}\bigr)$.
\item[\textup{(ii)}]
  Half-integer $m = k{+}\half$, $k \ge 0$:
  $\operatorname{val}(m)$ is an exact rational (Wallis product partial convergent).
\item[\textup{(iii)}]
  $\operatorname{val}(m) \sim \sqrt{4m/\pi}$
  as $m \to \infty$ by Stirling's formula.
\end{enumerate}
\end{corollary}

\begin{proof}
Part~\textup{(i)} is the integer case of Theorem~\ref{thm:gamma}.
For part~\textup{(ii)}, substitute $m=k{+}\half$ into the Gamma quotient and
use $\Gamma(k{+}1)=k!$. Part~\textup{(iii)} is the corresponding Stirling
asymptotic for the same quotient.
\end{proof}

%======================================================================
\section{Casoratian Series Identity for \texorpdfstring{$\pi/4$}{pi/4}}%
\label{sec:casoratian}
%======================================================================

We give a Casoratian derivation of the series identity for~$\pi/4$.
The telescoping argument is internal to the $m=1$ recurrence, while
the identification of the continued-fraction limit as~$4/\pi$ comes
independently from the Euler/Wallis analysis of the previous section.

\subsection{Convergence}

\begin{proposition}\label{prop:conv}
The CF~\eqref{eq:pcf} with
$\alpha(n) = -n(2n{-}3)$, $\beta(n) = 3n{+}1$ converges.
\end{proposition}

\begin{proof}
By the \'Sleszy\'nski--Pringsheim theorem,
the CF converges if
$|\alpha(n)| \leq \beta(n)\,\beta(n{-}1)$ for $n \geq 2$.
We have
$\beta(n)\,\beta(n{-}1) - |\alpha(n)|
= (3n{+}1)(3n{-}2) - n(2n{-}3) = 7n^2 - 2$,
which is strictly positive for $n \geq 1$.
\end{proof}

\subsection{Casoratian recursion}

Let $h_n = n^2{+}3n{+}1$ (the polynomial solution from
Theorem~\ref{thm:conv1}) and $g_n = q_n/\dff{2n{-}1}$,
both satisfying the reduced recurrence~\eqref{eq:ured}.

\begin{lemma}[Casoratian recursion]\label{lem:cas_rec}
The Casoratian $W_n = h_n\,g_{n-1} - h_{n-1}\,g_n$ satisfies
$W_n = \frac{n}{2n{-}1}\,W_{n-1}$ for $n \geq 2$.
\end{lemma}

\begin{proof}
Substituting the recurrence~\eqref{eq:ured} for $h_n$ and $g_n$:
\begin{align*}
  W_n
  &= \frac{(3n{+}1)h_{n-1} - n\,h_{n-2}}{2n{-}1}\,g_{n-1}
    - h_{n-1}\,\frac{(3n{+}1)g_{n-1} - n\,g_{n-2}}{2n{-}1} \\
  &= \frac{n}{2n{-}1}(h_{n-1}\,g_{n-2} - h_{n-2}\,g_{n-1})
   = \frac{n}{2n{-}1}\,W_{n-1}.\qedhere
\end{align*}
\end{proof}

\begin{theorem}[Casoratian closed form]\label{thm:cas}
For all $n \geq 1$,
$W_n = n!/\dff{2n{-}1}$.
\end{theorem}

\begin{proof}
Base: $W_1 = 5\cdot 1 - 1\cdot 4 = 1 = 1!/1!!$.
Induction:
$W_n = \frac{n}{2n{-}1}\cdot\frac{(n{-}1)!}{\dff{2n{-}3}}
= \frac{n!}{\dff{2n{-}1}}$.
\end{proof}

\subsection{Main series identity}

\begin{theorem}\label{thm:4overpi}
Let $S$ be the CF with $\alpha(n) = -n(2n{-}3)$,
$\beta(n) = 3n{+}1$.
\begin{enumerate}
\item[\textup{(i)}]
  $S$ converges, with $|S_n - S| = O(2^{-n}/n^{7/2})$.
\item[\textup{(ii)}]
  $p_n = \dff{2n{-}1}\,(n^2{+}3n{+}1)$.
\item[\textup{(iii)}]
  $S = 4/\pi$, equivalently
\begin{equation}\label{eq:series}
  \frac{\pi}{4}
  = 1 - \sum_{k=1}^{\infty}
    \frac{k!}{\dff{2k{-}1}\,(k^2{+}3k{+}1)(k^2{+}k{-}1)}
  = 1 - \sum_{k=1}^{\infty}
    \frac{2^k}{\binom{2k}{k}(k^2{+}3k{+}1)(k^2{+}k{-}1)}.
\end{equation}
\end{enumerate}
\end{theorem}

\begin{proof}[Proof of \textup{(iii)}]
Part~\textup{(ii)} is exactly the closed numerator formula~\eqref{eq:pn_m1}
from Theorem~\ref{thm:conv1}. Part~\textup{(i)} combines
Proposition~\ref{prop:conv} with the tail estimate proved immediately below.
It remains to establish \textup{(iii)}. The ratio $g_n/h_n$ telescopes:
\[
  \frac{g_n}{h_n} - \frac{g_{n-1}}{h_{n-1}}
  = \frac{-W_n}{h_n\,h_{n-1}}.
\]
Summing from $k = 1$ to~$n$ and using $g_0/h_0 = 1$:
\[
  \frac{g_n}{h_n}
  = 1 - \sum_{k=1}^{n}
    \frac{k!}{\dff{2k{-}1}\,(k^2{+}3k{+}1)(k^2{+}k{-}1)}.
\]
As $n\to\infty$, the convergents $p_n/q_n = h_n/g_n$ tend to the
continued-fraction value $S$ by Proposition~\ref{prop:conv}.
The value $S = 4/\pi$ is established independently by
Lemma~\ref{lem:shift} and Step~3 of the proof of
Theorem~\ref{thm:gamma}: the raising operator $L$ maps
$P^{(0)},Q^{(0)}$ to scalar multiples of $P^{(1)},Q^{(1)}$, so
$S^{(1)} = 2\,S^{(0)} = 4/\pi$ by Theorem~\ref{thm:pi0}, with no
appeal to the present section. Therefore $g_n/h_n \to \pi/4$, and
passing to the limit in the telescoping identity yields
\eqref{eq:series}. The alternative form uses
$k!/\dff{2k{-}1} = 2^k/\binom{2k}{k}$
(see~\cite[\S1.2.6]{knuth1997} for the double-factorial convention).
\end{proof}

\begin{proof}[Proof of convergence rate \textup{(i)}]
By Stirling, $\binom{2k}{k} \sim 4^k/\sqrt{\pi k}$, so the
general term satisfies
$t_k \sim \sqrt{\pi k}/(2^k k^4) = O(2^{-k}/k^{7/2})$.
Since the tail is dominated by its first term,
$|S_N - \pi/4| = O(2^{-N}/N^{7/2})$.
At $N = 35$ the error drops below~$10^{-16}$.
\end{proof}

%======================================================================
\section{Gauss CF Non-Membership}\label{sec:gauss}
%======================================================================

\subsection{Factorisation over \texorpdfstring{$\mathbb{Q}(\sqrt{5})$}{Q(sqrt5)}}

Set $\varphi = (1{+}\sqrt{5})/2$, $\psi = (1{-}\sqrt{5})/2$.
Since $(k{+}\varphi)(k{+}\psi) = k^2{+}k{-}1$
and $(k{+}2{+}\varphi)(k{+}2{+}\psi) = k^2{+}5k{+}5$,
the ratio of consecutive series terms in~\eqref{eq:series} is
\begin{equation}\label{eq:ratio_hyp}
  \frac{t_{k+1}}{t_k}
  = \frac{(k{+}1)(k{+}\varphi)(k{+}\psi)}
         {2(k{+}\half)(k{+}2{+}\varphi)(k{+}2{+}\psi)}.
\end{equation}

\subsection{Pochhammer obstruction}

A standard $\hyp{p}{q}{\cdots}{\cdots}{z}$ has term ratio equal to
a product of linear factors $(k+a_i)/(k+b_j)$ with
$a_i,b_j\in\mathbb{Q}$.
Here the irreducible quadratics $k^2{+}k{-}1$ and $k^2{+}5k{+}5$
(both of discriminant~$5$) force the parameters into
$\mathbb{Q}(\sqrt5)\setminus\mathbb{Q}$.

\subsection{Gauss CF non-membership}

\begin{proposition}\label{prop:gauss}
The CF does not arise from any ${}_2F_1$ ratio via the Gauss
continued-fraction theorem.
\end{proposition}

\begin{proof}
The Gauss CF for
$\hyp21{A{+}1,\,B}{C{+}1}{z}\big/\hyp21{A,B}{C}{z}$
has unit-denominator coefficients $t_k$ determined by
$(A,B,C)\in\mathbb{Q}^3$.
Solving $t_1=c_1$, $t_2=c_2$, $t_3=c_3$ (where the $c_k$ are the
equivalence-transformed coefficients of our CF) yields a degree-6
polynomial system with six solutions; four are complex and two are
real with $C \approx -3.259$.
For these,
$|t_4 - c_4| \approx 5.9$, so no match exists.
\end{proof}

\subsection{\texorpdfstring{${}_2F_3$}{2F3} identification}

\begin{remark}\label{rem:2f3}
The term ratio~\eqref{eq:ratio_hyp} identifies the series as
\[
  {}_2F_3\!\left(\varphi,\psi;\;\half,\;2{+}\varphi,\;2{+}\psi;\;
  \half\right)
\]
over $\mathbb{Q}(\sqrt5)$.
Since both the parameters and the function class depend on the
extension $\mathbb{Q}(\sqrt5)/\mathbb{Q}$, the series does not
reduce to any standard ${}_pF_q$ with rational parameters.
\end{remark}

%======================================================================
\section{Euler CF Duality}\label{sec:euler}
%======================================================================

The $m=0$ CF decomposes as
$S^{(0)} = \beta(0) + \alpha(1)/T = 1 - 1/T$
where~$T$ is the inner tail starting at $n = 1$.
Since $S^{(0)} = 2/\pi$ (Theorem~\ref{thm:pi0}):

\begin{proposition}\label{prop:euler}
$T = \pi/(\pi - 2)$, equivalently
$T/(T{-}1) = \pi/2$.
\end{proposition}

\begin{proof}
$T = 1/(1 - S^{(0)}) = 1/(1 - 2/\pi) = \pi/(\pi - 2)$.
The identity $T/(T{-}1) = \pi/2$ then follows from
$T - 1 = (2)/(\pi - 2)$.
\end{proof}

\begin{remark}
This connects our PCF to the classical identity~\eqref{eq:pi-id}, which comes
from $\arcsin(1)$~\cite[\S4.24.2]{dlmf}.
The partial numerators of the $m=0$ and $m=1$ families differ
by a linear term:
$\alpha_0(n) - \alpha_1(n) = -n(2n{-}1) + n(2n{-}3) = -2n$.
\end{remark}

%======================================================================
\section{Comparison with Brouncker's CF}\label{sec:brouncker}
%======================================================================

Lord Brouncker's classical CF for $4/\pi$ is
\[
  \frac{4}{\pi}
  = 1 + \cfrac{1^2}{2 + \cfrac{3^2}{2 + \cfrac{5^2}{2 + \cdots}}},
\]
with $\alpha^B(n) = (2n{-}1)^2$ and constant $\beta^B(n) = 2$.
Both CFs converge to~$4/\pi$ but differ structurally:

\begin{center}
\begin{tabular}{@{}lcc@{}}
  \toprule
  & \textbf{Brouncker} & \textbf{Novel ($m{=}1$)} \\
  \midrule
  $\alpha(n)$ & $(2n{-}1)^2$ & $-n(2n{-}3)$ \\
  $\beta(n)$ & $2$ (constant) & $3n{+}1$ (linear) \\
  $p_n$ & $(2n{+}1)!!$ & $\dff{2n{-}1}(n^2{+}3n{+}1)$ \\
  Series & Leibniz $\sum (-1)^k/(2k{+}1)$ &
    Eq.~\eqref{eq:series} \\
  Rate & $O(1/n)$ & $O(2^{-n})$ \\
  \bottomrule
\end{tabular}
\end{center}

\begin{center}
\small
\begin{tabular}{@{}rrrcrrc@{}}
  \toprule
  & \multicolumn{3}{c}{\textbf{Brouncker}} &
    \multicolumn{3}{c}{\textbf{Novel}} \\
  \cmidrule(lr){2-4}\cmidrule(lr){5-7}
  $n$ & $p_n$ & $q_n$ & $S_n$ & $p_n$ & $q_n$ & $S_n$ \\
  \midrule
  0 & 1 & 1 & 1.000000 & 1 & 1 & 1.000000 \\
  1 & 3 & 2 & 1.500000 & 5 & 4 & 1.250000 \\
  2 & 15 & 13 & 1.153846 & 33 & 26 & 1.269231 \\
  3 & 105 & 76 & 1.381579 & 285 & 224 & 1.272321 \\
  4 & 945 & 789 & 1.197719 & 3045 & 2392 & 1.272993 \\
  5 & 10395 & 7734 & 1.344065 & 38745 & 30432 & 1.273166 \\
  8 & 34459425 & 28018665 & 1.229874 &
    180405225 & 141690240 & 1.273237 \\
  \bottomrule
  \multicolumn{7}{r}{$4/\pi \approx 1.273\,239\,545$}
\end{tabular}
\end{center}

\noindent
At $n = 8$, Brouncker's error is $\sim\!4\times 10^{-2}$ while the
novel CF error is $\sim\!3\times 10^{-6}$, a factor of~$10^4$.
The two CFs are not related by any equivalence transformation,
contraction, or M\"obius substitution.

%======================================================================
\section{Convergence Theory}\label{sec:conv}
%======================================================================

\subsection{Worpitzky condition}

\begin{lemma}\label{lem:worpitzky}
For both families, the normalised partial numerators
$|\tilde\alpha_n| = |\alpha_n|/(\beta_n\,\beta_{n-1})$ satisfy
$|\tilde\alpha_n| < 1/4$ for all sufficiently large~$n$.
\end{lemma}

\begin{proof}
\textbf{Log family.}
$|\tilde\alpha_n| \to k/(k{+}1)^2 < 1/4$ for $k > 1$
(since $(k{-}1)^2 > 0$).

\textbf{Pi family.}
$|\tilde\alpha_n| = n(2n{-}1)/((3n{+}1)(3n{-}2)) \to 2/9 < 1/4$.
At $n = 1$: $|\tilde\alpha_1| = 1/4$ (borderline);
strict inequality holds for $n \ge 2$, and convergence at
equality at finitely many~$n$ follows from the classical
extension~\cite[\S4.2]{lorentzen2008}.
\end{proof}

\subsection{Bracketing}

\begin{lemma}\label{lem:bracket}
If $\alpha_n < 0$ and $\beta_n > 0$ for all $n \ge 1$, and the CF
converges to~$L$, then
$\min(C_N,C_{N+1}) \le L \le \max(C_N,C_{N+1})$.
\end{lemma}

\begin{proof}
The determinant $p_n q_{n-1} - p_{n-1} q_n
= \prod_{j=1}^n (-\alpha_j) > 0$
shows $C_n - C_{n-1}$ alternates in sign.
See~\cite[Theorem~4.35]{lorentzen2008}.
\end{proof}

%======================================================================
\section{Meta-Family}\label{sec:meta}
%======================================================================

The unified template
$\alpha(n) = -(\alpha_0 n^2 + \beta_0 n)$,
$\beta(n) = \gamma_0 n + \delta_0$
contains both the Logarithmic Ladder
($\beta_0 = 0$, $\delta_0 = \alpha_0$, $\gamma_0 = \alpha_0 + 1$)
and the Pi Family
($\alpha_0 = 2$, $\gamma_0 = 3$, $\delta_0 = 1$)
as special cases.
Two additional sub-families were found and checked at $100$ decimal digits
with forward depth $N=2000$:
\begin{itemize}
  \item $(\alpha_0, \beta_0, \gamma_0, \delta_0) = (1,0,4,2)
    \to 2/\ln 3$;
  \item $(4,0,8,4) \to 4/\ln 3$.
\end{itemize}
Both are instances of the Logarithmic Ladder at~$k = 3$
(and its equivalence-transformed double).

%======================================================================
\section{Quadratic GCF Constants}\label{sec:vquad}
%======================================================================

Define $V(A,B,C) = 1 + \mathbf{K}_{n \ge 1}\; 1/(An^2 + Bn + C)$.

\begin{proposition}\label{prop:irrat}
For all $(A,B,C)$ with $A \ge 1$, $C \ge 1$, and
$An^2 + Bn + C > 0$ for $n \ge 1$, $V(A,B,C)$ is irrational.
\end{proposition}

\begin{proof}
The convergent denominators satisfy
$Q_n = (An^2{+}Bn{+}C)\,Q_{n-1} + Q_{n-2}$
with $\log Q_n \gtrsim n \log n$.
The Wronskian $|P_n Q_{n-1} - P_{n-1} Q_n| = 1$ for all~$n$.
By Euler's irrationality criterion, the limit is irrational.
\end{proof}

Our systematic scan found \textbf{482 distinct} such constants.
All values in this paragraph were first computed at $300$ decimal digits with
forward depth $N=2000$, and the final Arb enclosures are recorded in
Appendix~\ref{sec:cert}.
A PSLQ~\cite{ferguson1999} integer relation search at 300-digit precision over 15
standard constants
$\{\pi, \pi^2, e, \ln 2, \ln 3, \gamma, G, \zeta(3), \zeta(5),
\sqrt{2}, \sqrt{3}, \sqrt{5}, \varphi, \sqrt{\pi}, \pi\ln 2\}$
finds no relation for any $V$-constant, and none is algebraic of
degree~$\le 8$ with coefficients~$\le 10^5$.
Pairwise cross-relation tests on 91 representative constants
(435~pairs, PSLQ with coefficients up to~500) and triple-relation
tests (455~triples, coefficients up to~200) all return negative.
A targeted sweep against 588 expressions built from Gamma ratios
and Pi Family values also produces no matches.

These computations have a mathematical role beyond data collection.
Proposition~\ref{prop:irrat} gives a uniform irrationality theorem for the
quadratic family, and the negative PSLQ evidence indicates that the 482 values
are not simple reparameterisations of the standard constants already in the
search basis.

The only structural relation found is the \emph{shift symmetry}
$V(1,B,C) - V(1,-B,C) \in \mathbb{Z}$ for all tested $B, C$,
following from the palindromic structure of the GCF denominators.
This gives strong (but non-rigorous) evidence of mutual algebraic
independence.

\begin{remark}
An earlier PSLQ search returned $[0, 1, -2, 1]$ relative to
$\{V, \sqrt{5}, \varphi, 1\}$; this is the trivial identity
$\sqrt{5} - 2\varphi + 1 = 0$ with zero coefficient on~$V$.
\end{remark}

%======================================================================
\appendix
\section{Numerical Certification}\label{sec:cert}
%======================================================================

All results are certified using Arb ball arithmetic
(\texttt{python-flint}~0.8.0, $10\,000$-bit precision)
at depth $N = 5000$.
By Lemma~\ref{lem:bracket}, the limit lies in
$[\min(C_N, C_{N-1}),\, \max(C_N, C_{N-1})]$.
Table~\ref{tab:cert} summarises the main certified enclosures.
These computations serve as certification of the proved formulas and of the
catalogue values recorded in Section~\ref{sec:vquad}.

\begin{table}[ht]
\centering
\caption{Arb-certified intervals (depth~5000)}\label{tab:cert}
\begin{tabular}{@{}lccc@{}}
\toprule
PCF & Target & Bracket width & Digits \\
\midrule
$\alpha{=}{-}2n^2,\;\beta{=}3n{+}2$
  & $1/\ln 2$    & $<10^{-1509}$ & 1509 \\
$\alpha{=}{-}3n^2,\;\beta{=}4n{+}3$
  & $1/\ln(3/2)$ & $<10^{-2265}$ & 2265 \\
$\alpha{=}{-}n(2n{-}1),\;\beta{=}3n{+}1$
  & $2/\pi$      & $<10^{-1508}$ & 1508 \\
$\alpha{=}{-}n(2n{-}3),\;\beta{=}3n{+}1$
  & $4/\pi$      & $<10^{-1518}$ & 1518 \\
\bottomrule
\end{tabular}
\end{table}

For the $m = 1$ CF, backward recurrence at depth~500 with
80-digit arithmetic confirms
\[
  S = 1.27323954473516268615107010698011489627567716592365158998\ldots
\]
matching $4/\pi$ to all 60 digits.

%======================================================================
\section{Discovery Methodology}\label{sec:method}
%======================================================================

\begin{enumerate}
\item \textbf{MITM sweep.}
  Fix $\beta(n) = dn{+}e$; enumerate
  $\alpha(n) = pn^2{+}qn{+}r$ with $|p|,|q|,|r| \le 5$;
  match PCF values against a database of 200 known constants.
\item \textbf{Modular signature sweep.}
  Vary $d \in \{2,\ldots,8\}$, $e \in \{1,\ldots,8\}$.
  The $d{=}3$ family produced both $\pi$- and $\ln$-hits.
\item \textbf{Parametric generalisation.}
  Sweep $c$ in $\alpha(n) = -n(2n{-}c)$ and $k$ in
  $\alpha(n) = -kn^2$, discovering the full infinite families.
\end{enumerate}

%======================================================================
\section{Open Questions}\label{sec:open}
%======================================================================

\begin{enumerate}
\item \textbf{Prove Conjecture~\ref{conj:ratio}}
  (the ratio formula for general~$m$).
  A proof would complete Theorem~\ref{thm:gamma} unconditionally.
  Clean closed forms for~$p_n^{(m)}$ and~$q_n^{(m)}$ for
  general~$m > 0$ would likely suffice.
\item \textbf{Irrationality measures.}
  Do the convergent sequences yield new irrationality measures
  for $\ln 2$, $\ln(3/2)$, or~$1/\pi$ via Ap\'ery-style arguments?
\item \textbf{Higher denominators.}
  Do analogous families exist for $d \ge 5$, producing $\zeta(3)$,
  Catalan's constant, or polylogarithms?
\item \textbf{Quadratic GCFs.}
  Determine whether
  $V(3,1,1) \approx 1.19737399068835760\ldots$
  has any closed form.
  This constant resists $3\,400{+}$ PSLQ tests at $300$ decimal digits,
  using continued-fraction values computed to depth $N=2000$, against
  15 special-function families~\cite{raayoni2021}.
\item \textbf{Algebraic independence.}
  Prove (or disprove) that the 482 quadratic GCF constants are
  mutually algebraically independent.
\end{enumerate}

%======================================================================
\section{Exact Convergents}\label{app:conv}
%======================================================================

\textbf{Log ($k{=}2$):} $p_n = (n{+}1)!\cdot 2^{n+1}$.

\smallskip
\begin{center}
\begin{tabular}{@{}rrrl@{}}
\toprule
$n$ & $p_n$ & $q_n$ & $C_n$ \\
\midrule
0 & 2 & 1 & 2.000000 \\
1 & 8 & 5 & 1.600000 \\
2 & 48 & 32 & 1.500000 \\
3 & 384 & 262 & 1.465649 \\
4 & 3840 & 2644 & 1.452345 \\
5 & 46080 & 31832 & 1.447584 \\
\bottomrule
\end{tabular}
\end{center}
\smallskip\noindent Limit: $1/\ln 2 = 1.442695\ldots$

\medskip
\textbf{Pi ($m{=}0$):} $p_n = (2n{+}1)!!$.

\smallskip
\begin{center}
\begin{tabular}{@{}rrrl@{}}
\toprule
$n$ & $p_n$ & $q_n$ & $C_n$ \\
\midrule
0 & 1 & 1 & 1.000000 \\
1 & 3 & 4 & 0.750000 \\
2 & 15 & 22 & 0.681818 \\
3 & 105 & 160 & 0.656250 \\
4 & 945 & 1464 & 0.645492 \\
5 & 10395 & 16224 & 0.640688 \\
\bottomrule
\end{tabular}
\end{center}
\smallskip\noindent Limit: $2/\pi = 0.636620\ldots$

\medskip
\textbf{Pi ($m{=}1$):}
$p_n = \dff{2n{-}1}\,(n^2{+}3n{+}1)$.

\smallskip
\begin{center}
\begin{tabular}{@{}rrl@{}}
\toprule
$n$ & $p_n$ & $q_n$ \\
\midrule
0 & 1 & 1 \\
1 & 5 & 4 \\
2 & 33 & 26 \\
3 & 285 & 224 \\
4 & 3\,045 & 2\,392 \\
5 & 38\,745 & 30\,432 \\
6 & 571\,725 & 449\,040 \\
7 & 9\,594\,585 & 7\,535\,616 \\
8 & 180\,405\,225 & 141\,690\,240 \\
\bottomrule
\end{tabular}
\end{center}
\smallskip\noindent Limit: $4/\pi = 1.273239\ldots$

%======================================================================
\section{Reproducibility}\label{app:repro}
%======================================================================

Python~3.12, \texttt{mpmath}~1.4.1,
\texttt{python-flint}~0.8.0 (Arb),
\texttt{sympy}~1.14.0.
Forward recurrence~\eqref{eq:rec} with
$p_{-1}{=}1$, $p_0{=}\beta(0)$, $q_{-1}{=}0$, $q_0{=}1$.
Runtime: ${<}\,2$~hours on a standard desktop.

%======================================================================
\section*{Acknowledgements}
%======================================================================

Inspired by the Ramanujan Machine~\cite{raayoni2021,elimelech2024}.
Computations performed with mpmath and python-flint (Arb).

%======================================================================
\begin{thebibliography}{12}
%======================================================================

\bibitem{apery1979}
R.~Ap\'ery,
\textit{Irrationalit\'e de $\zeta(2)$ et $\zeta(3)$},
Ast\'erisque \textbf{61} (1979), 11--13.

\bibitem{cuyt2008}
A.~Cuyt, V.\,B.~Petersen, B.~Verdonk, H.~Waadeland,
and W.\,B.~Jones,
\emph{Handbook of Continued Fractions for Special Functions},
Springer, 2008.

\bibitem{dlmf}
F.\,W.\,J.~Olver, A.\,B.~Olde~Daalhuis, D.\,W.~Lozier,
et~al.\ (eds.),
\emph{NIST Digital Library of Mathematical Functions},
Release~1.2.3, \url{https://dlmf.nist.gov/}, 2024.

\bibitem{elimelech2024}
D.~Elimelech et~al.,
\textit{Algorithm-assisted discovery of an intrinsic order among
mathematical constants},
Proc.\ Natl.\ Acad.\ Sci.\ USA \textbf{121}(25) (2024),
e2321440121.

\bibitem{ferguson1999}
H.\,R.\,P.~Ferguson, D.\,H.~Bailey, and S.~Arwade,
\textit{Analysis of PSLQ, an integer relation finding algorithm},
Math.\ Comp.\ \textbf{68} (1999), 351--369.

\bibitem{knuth1997}
D.\,E.~Knuth,
\emph{The Art of Computer Programming},
vol.~1, 3rd ed., Addison-Wesley, 1997.

\bibitem{lorentzen2008}
L.~Lorentzen and H.~Waadeland,
\emph{Continued Fractions with Applications},
2nd ed., Atlantis Press, 2008.

\bibitem{oeis_A028387}
OEIS Foundation,
\textit{A028387: $n(n{+}2)+1$},
The On-Line Encyclopedia of Integer Sequences,
\url{https://oeis.org/A028387}.

\bibitem{raayoni2021}
G.~Raayoni, S.~Gottlieb, Y.~Manor, G.~Pisha, Y.~Harris,
U.~Mendlovic, D.~Haviv, Y.~Hadad, and I.~Kaminer,
\textit{Generating conjectures on fundamental constants with the
Ramanujan Machine},
Nature \textbf{590} (2021), 67--73.

\bibitem{wall1948}
H.\,S.~Wall,
\emph{Analytic Theory of Continued Fractions},
Van Nostrand, 1948; reprinted Chelsea, 1967.

\end{thebibliography}

\end{document}
